\section{Introduction}
\label{sec:introduction}

Concrete compressive strength is a central response variable in mixture design, quality control, and structural-material assessment. Public datasets have made this problem a common example for comparing statistical and machine-learning models because the predictors are familiar mixture descriptors, the response has direct engineering meaning, and the modelling task can be explained without a large domain-specific preprocessing pipeline. The UCI Concrete Compressive Strength dataset, attributed to Yeh, provides one such public benchmark and records mixture constituents, curing age, and compressive strength for 1,030 observations in its official repository description \citep{Yeh1998_UCIConcreteDataset}. Yeh's earlier modelling study also established concrete-strength prediction from mixture variables and curing age as a useful applied test case for data-driven models \citep{Yeh1998_ConcreteANN}.

Recent concrete-strength studies include more flexible computational-intelligence methods and broader model comparisons \citep{NunezMaraniFlahNehdi2021_ConcreteMLReview}. This manuscript intentionally takes a narrower route. Rather than claiming a state-of-the-art concrete predictor, it uses a simple regularized linear model to demonstrate how an auditable manuscript can keep methods, tables, figures, provenance, and claim limits aligned. That focus is deliberate: an interpretable baseline is often easier to inspect, reproduce, and challenge than a high-capacity model.

Ridge regression adds an \(L_2\) penalty to ordinary least-squares estimation and is commonly used when predictors are correlated or when coefficient shrinkage is desirable \citep{HoerlKennard1970_Ridge}. Cross-validation provides a resampling-based approach for model assessment and model-selection support \citep{Stone1974_CrossValidation}. These methods are mature; the contribution of this example is not methodological novelty, but disciplined reporting under constrained evidence.

A critical provenance limitation applies throughout the manuscript. The current generated artifacts are labelled \texttt{DETERMINISTIC\_FALLBACK\_NOT\_UCI}. Therefore, the UCI dataset is cited only as the intended public-data context and schema reference. Numerical results in this manuscript describe the local fallback artifacts generated for the demonstration pipeline, not verified empirical results from the official UCI repository. The objectives are to: (i) document the local fallback dataset and ridge-regression workflow; (ii) report held-out test performance, residual patterns, coefficient summaries, and ablation results using the generated tables and figures; and (iii) state the claim boundaries needed before any final registry or citation-mapping pass.
